% GENERATED FILE. Edit canonical lessons, then run npm run generate:books.
% canonical-source-hash: fnv1a64:ea92368f4f3f4211
% canonical-lessons: JA-C326-kisoku, JA-C326-yaburu, JA-C326-kinshi, JA-C326-ihan, JA-C326-batsu

\chapter{Rules}
\label{ch:ja-rules}
\hlchaptermodality{326}

\begin{chapteropening}
\textbf{By the end of this chapter:} \emph{I can say what a rule is and why it is there.}

Complete the last lesson of chapter 326: I can say what a rule is and why it is there.
\end{chapteropening}

\section[\texorpdfstring{\ja{きそく} (kisoku)}{きそく (kisoku)}]{\ja{きそく} (kisoku) — a rule; a regulation}

\label{lesson:JA-C326-kisoku}



\begin{quote}
Before the new one: say the Japanese for after all, then the Japanese for freedom.
\end{quote}

\subsection*{You'll want to know: \ja{きそく}}
\textbf{\ja{きそく}} — \emph{kisoku} — \textquotedblleft{}a rule; a regulation\textquotedblright{}.

\begin{quote}
\textbf{\ja{あぶないから、この} \ja{きそくが} \ja{あります。}} — \emph{abunai kara, kono kisoku ga arimasu} — \textquotedblleft{}This rule is there because it is dangerous.\textquotedblright{}
\end{quote}

\begin{cousinweb}
\textbf{\ja{きそく}} (\emph{kisoku}) is Sino-Japanese (kango): two Chinese characters read with their on'yomi, readings borrowed from Chinese. They are \emph{ki}, \textquotedblleft{}standard\textquotedblright{} (Middle Chinese \emph{kjwie}) and \emph{soku}, \textquotedblleft{}rule\textquotedblright{} (\emph{tsok}) (KANJIDIC2; Middle Chinese in Baxter's transcription).
\end{cousinweb}

\begin{grammarlens}[title={What you've built}]
One more word for giving reasons and opinions.
\end{grammarlens}

\subsection*{Your turn}
\begin{itemize}
\raggedright
  \item \emph{Say it:} \emph{kisoku}
  \item \emph{Say it:} \emph{kisoku}, once more
  \item \emph{Say it:} \emph{abunai kara, kono kisoku ga arimasu} — This rule is there because it is dangerous.
  \item \emph{Recall:} say the Japanese for after all, then the Japanese for freedom, then say \emph{kisoku} again
\end{itemize}

\begin{tcolorbox}[breakable,colback=teal!4,colframe=teal!35!black,title={Before you move on}]
What does \ja{きそく} mean? (\textquotedblleft{}A rule; a regulation\textquotedblright{}.) Say it once more.
\end{tcolorbox}
\section[\texorpdfstring{\ja{やぶる} (yaburu)}{やぶる (yaburu)}]{\ja{やぶる} (yaburu) — to break a rule; to break (a promise); to tear}

\label{lesson:JA-C326-yaburu}



\begin{quote}
Before the new one: say the Japanese for freedom, then the Japanese for a regulation.
\end{quote}

\subsection*{You'll want to know: \ja{やぶる}}
\textbf{\ja{やぶる}} — \emph{yaburu} — \textquotedblleft{}to break a rule; to break (a promise); to tear\textquotedblright{}.

\begin{quote}
\textbf{\ja{やくそくを} \ja{やぶったから、ともだちが} \ja{おこりました。}} — \emph{yakusoku o yabutta kara, tomodachi ga okorimashita} — \textquotedblleft{}I broke my promise, so my friend got angry.\textquotedblright{}
\end{quote}

Polite \textbf{\ja{やぶります}}, \emph{yaburimasu}; a godan verb. Past (polite): \textbf{\ja{かれが} \ja{やくそくを} \ja{やぶりました。}}, \emph{kare ga yakusoku o yaburimashita}.

\begin{cousinweb}
\textbf{\ja{やぶる}} (\emph{yaburu}) is a native Japanese word (wago), inherited rather than borrowed from Chinese, and its older origin is not settled. Its kanji was attached to it later as a kun'yomi, a native reading given to a Chinese character (JMdict, KANJIDIC2).
\end{cousinweb}

\begin{grammarlens}[title={What you've built}]
One more word for giving reasons and opinions.
\end{grammarlens}

\subsection*{Your turn}
\begin{itemize}
\raggedright
  \item \emph{Say it:} \emph{yaburu}
  \item \emph{Say it:} \emph{yaburu}, once more
  \item \emph{Say it:} \emph{yakusoku o yabutta kara, tomodachi ga okorimashita} — I broke my promise, so my friend got angry.
  \item \emph{Recall:} say the Japanese for freedom, then the Japanese for a regulation, then say \emph{yaburu} again
\end{itemize}

\begin{tcolorbox}[breakable,colback=teal!4,colframe=teal!35!black,title={Before you move on}]
What does \ja{やぶる} mean? (\textquotedblleft{}To break a rule; to break (a promise); to tear\textquotedblright{}.) Say it once more.
\end{tcolorbox}
\section[\texorpdfstring{\ja{きんし} (kinshi)}{きんし (kinshi)}]{\ja{きんし} (kinshi) — a ban; a prohibition}

\label{lesson:JA-C326-kinshi}



\begin{quote}
Before the new one: say the Japanese for a regulation, then the Japanese for to break a rule.
\end{quote}

\subsection*{You'll want to know: \ja{きんし}}
\textbf{\ja{きんし}} — \emph{kinshi} — \textquotedblleft{}a ban; a prohibition\textquotedblright{}.

\begin{quote}
\textbf{\ja{かわは} \ja{あぶないから、きんしに} \ja{なりました。}} — \emph{kawa wa abunai kara, kinshi ni narimashita} — \textquotedblleft{}The river is dangerous, so it has been put off limits.\textquotedblright{}
\end{quote}

\begin{cousinweb}
\textbf{\ja{きんし}} (\emph{kinshi}) is Sino-Japanese (kango): two Chinese characters read with their on'yomi, readings borrowed from Chinese. They are \emph{kin}, \textquotedblleft{}prohibition\textquotedblright{} (Middle Chinese \emph{kimH}) and \emph{shi}, \textquotedblleft{}stop\textquotedblright{} (\emph{tsyiX}) (KANJIDIC2; Middle Chinese in Baxter's transcription).
\end{cousinweb}

\begin{grammarlens}[title={What you've built}]
One more word for giving reasons and opinions.
\end{grammarlens}

\subsection*{Your turn}
\begin{itemize}
\raggedright
  \item \emph{Say it:} \emph{kinshi}
  \item \emph{Say it:} \emph{kinshi}, once more
  \item \emph{Say it:} \emph{kawa wa abunai kara, kinshi ni narimashita} — The river is dangerous, so it has been put off limits.
  \item \emph{Recall:} say the Japanese for a regulation, then the Japanese for to break a rule, then say \emph{kinshi} again
\end{itemize}

\begin{tcolorbox}[breakable,colback=teal!4,colframe=teal!35!black,title={Before you move on}]
What does \ja{きんし} mean? (\textquotedblleft{}A ban; a prohibition\textquotedblright{}.) Say it once more.
\end{tcolorbox}
\section[\texorpdfstring{\ja{いはん} (ihan)}{いはん (ihan)}]{\ja{いはん} (ihan) — a violation; breaking the rules}

\label{lesson:JA-C326-ihan}



\begin{quote}
Before the new one: say the Japanese for to break a rule, then the Japanese for a ban.
\end{quote}

\subsection*{You'll want to know: \ja{いはん}}
\textbf{\ja{いはん}} — \emph{ihan} — \textquotedblleft{}a violation; breaking the rules\textquotedblright{}.

\begin{quote}
\textbf{\ja{ここに} \ja{くるまを} \ja{とめたら、いはんです。}} — \emph{koko ni kuruma o tometara, ihan desu} — \textquotedblleft{}Parking your car here is against the rules.\textquotedblright{}
\end{quote}

\begin{cousinweb}
\textbf{\ja{いはん}} (\emph{ihan}) is Sino-Japanese (kango): two Chinese characters read with their on'yomi, readings borrowed from Chinese. They are \emph{i}, \textquotedblleft{}differ\textquotedblright{} (Middle Chinese \emph{hjw+j}) and \emph{han}, \textquotedblleft{}anti-\textquotedblright{} (\emph{pjonX}) (KANJIDIC2; Middle Chinese in Baxter's transcription).
\end{cousinweb}

\begin{grammarlens}[title={What you've built}]
One more word for giving reasons and opinions.
\end{grammarlens}

\subsection*{Your turn}
\begin{itemize}
\raggedright
  \item \emph{Say it:} \emph{ihan}
  \item \emph{Say it:} \emph{ihan}, once more
  \item \emph{Say it:} \emph{koko ni kuruma o tometara, ihan desu} — Parking your car here is against the rules.
  \item \emph{Recall:} say the Japanese for to break a rule, then the Japanese for a ban, then say \emph{ihan} again
\end{itemize}

\begin{tcolorbox}[breakable,colback=teal!4,colframe=teal!35!black,title={Before you move on}]
What does \ja{いはん} mean? (\textquotedblleft{}A violation; breaking the rules\textquotedblright{}.) Say it once more.
\end{tcolorbox}
\section[\texorpdfstring{\ja{ばつ} (batsu)}{ばつ (batsu)}]{\ja{ばつ} (batsu) — a punishment; a penalty}

\label{lesson:JA-C326-batsu}



\begin{quote}
Before the new one: say the Japanese for a ban, then the Japanese for a violation.
\end{quote}

\subsection*{You'll want to know: \ja{ばつ}}
\textbf{\ja{ばつ}} — \emph{batsu} — \textquotedblleft{}a punishment; a penalty\textquotedblright{}.

\begin{quote}
\textbf{\ja{きそくを} \ja{やぶったので、ばつが} \ja{ありました。}} — \emph{kisoku o yabutta node, batsu ga arimashita} — \textquotedblleft{}He broke the rules, so there was a punishment.\textquotedblright{}
\end{quote}

\begin{cousinweb}
\textbf{\ja{ばつ}} (\emph{batsu}) is Sino-Japanese (kango): one Chinese character read with its on'yomi, reading borrowed from Chinese. It is \emph{batsu}, \textquotedblleft{}punishment\textquotedblright{} (Middle Chinese \emph{bjot}) (KANJIDIC2; Middle Chinese in Baxter's transcription).
\end{cousinweb}

\begin{grammarlens}[title={What you've built}]
One more word for giving reasons and opinions.
\end{grammarlens}

\subsection*{Your turn}
\begin{itemize}
\raggedright
  \item \emph{Say it:} \emph{batsu}
  \item \emph{Say it:} \emph{batsu}, once more
  \item \emph{Say it:} \emph{kisoku o yabutta node, batsu ga arimashita} — He broke the rules, so there was a punishment.
  \item \emph{Recall:} say the Japanese for a ban, then the Japanese for a violation, then say \emph{batsu} again
\end{itemize}

\begin{tcolorbox}[breakable,colback=teal!4,colframe=teal!35!black,title={Before you move on}]
What does \ja{ばつ} mean? (\textquotedblleft{}A punishment; a penalty\textquotedblright{}.) Say it once more.
\end{tcolorbox}
